\begin{table}[ht]
\centering\small
\begin{tabular}{p{0.75\linewidth}r}
\toprule
bukti & nilai \\
\midrule
max\_abs\_spearman\_n\_vs\_perilaku & 0.010 \\
max\_MI\_berlebih\_vs\_acuan\_acak & 0.025 \\
silhouette\_cluster\_n\_saja & 0.670 \\
ARI(cluster\_n, cluster\_perilaku) & -0.000 \\
ARI(cluster\_semua\_kolom, cluster\_n) & 0.083 \\
ARI(cluster\_semua\_kolom, cluster\_perilaku) & 0.398 \\
jumlah\_kolom\_n & 15.000 \\
ARI(cluster\_n, kota\_tier) & -0.000 \\
ARI(cluster\_n, perangkat) & -0.000 \\
\bottomrule
\end{tabular}
\caption{Bukti bahwa n01–n15 membentuk struktur sendiri yang tidak terkait perilaku/profil pelanggan.}
\label{tab:t3_bukti_n}
\end{table}
